\documentclass[fleqn]{article}
\usepackage[spanish,es-noshorthands]{babel}
\usepackage[utf8]{inputenc} 
\usepackage[papersize={6.5in,8.5in},left=1cm, right=1cm, top=1.5cm, bottom=1.7cm]{geometry}
\usepackage{mathexam}
\usepackage{amsmath}
\usepackage{graphicx}

\ExamClass{\includegraphics[height=16pt]{Images/logo-sed.png} Álgebra $8^{\circ}$}
\ExamName{Expresiones algebraicas}
\ExamHead{\includegraphics[height=16pt]{Images/logo-colegio.png} IEDAB}
\newcommand{\LineaNombre}{%
\par
\vspace{\baselineskip}
Nombre:\hrulefill \; Curso: \underline{\hspace*{48pt}} \; Fecha: \underline{\hspace*{2.5cm}} \relax
\par}
\let\ds\displaystyle

\begin{document}
\ExamInstrBox{
Respuesta sin justificar mediante procedimiento no será tenida en cuenta en la calificación. Escriba sus respuestas en el espacio indicado. Tiene 45 minutos para contestar esta prueba.}
\LineaNombre
\begin{enumerate}
   \item Exprese en lenguaje algebraico las siguientes expresiones:
   \begin{enumerate}
    \item El triple de un número:
    \item El triple de la diferencia de dos números:
    \item La suma del triple de un número y el cuádruple de otro:
    \item La raíz cuadrada de la diferencia de dos números:
    \item La diferencia entre el triple de un número y el quíntuple de otro:
   \end{enumerate}
  \item Traduce al lenguaje verbal las siguientes expresiones algebraicas:
  \begin{enumerate}
   \item $\dfrac{b}{3}$\noanswer[15pt]
   \item $(x-2y)^{2}$\noanswer[15pt]
   \item $3c+5$\noanswer
   \item $2n^{2}-3m$\noanswer
  \end{enumerate}
  \item Complete la siguiente tabla, calculando el valor numérico de la expresión dada:
  {%
\newcommand{\mc}[3]{\multicolumn{#1}{#2}{#3}}
\begin{center}
\begin{tabular}{|l|l|l|l|l|}\hline
\mc{1}{|c|}{\textbf{x}} & \mc{1}{c|}{\textbf{y}} & \mc{1}{c|}{\textbf{z}} & \mc{1}{c|}{\textbf{Expresión alg.}} & \mc{1}{c|}{\textbf{Valor numérico}}\\\hline
2 & 3 & 5 & $x+y-z$ & \\\hline
$-2$ & 3 & $-4$ & $(x+y)^{2}-z$ & \\\hline
$-3$ & $-2$ & $-1$ & $2x^{2}-3y+4z$ & \\\hline
  \end{tabular}
  \end{center}
}%
\noanswer
\newpage
  \item Si $A=(x+4)$ y $B=(2x-3)$, halle lo siguiente:
  \begin{enumerate}
   \item $A+B=$\noanswer[25pt]
   \item $A-B=$\noanswer[25pt]
   \item $2\cdot A=$\noanswer[25pt]
   \item $-5B=$\noanswer[25pt]
  \end{enumerate}
  \item Resuelve los siguientes problemas, planteando la ecuación y resolviéndola. Compruebe sus resultados
  \begin{enumerate}
   \item ¿Qué número aumentado en 32 da 59?\noanswer
   \item Entre Antonio y Mario tienen 970 dólares. Sabiendo que Antonio tiene 75 dólares más que Mario, calcular los dólares que tiene cada uno\noanswer
   \item Repartir 688 dólares entre Juan, Pedro y Kevin, sabiendo que Juan debe recibir 38 dólares más que Pedro y 12 dólares menos que Kevin.
  \end{enumerate}
\end{enumerate}
\end{document}
